\documentclass[11pt,twocolumn]{article}
\usepackage[margin=1in,top=1in]{geometry}
\usepackage{times}
\usepackage{latexsym}
\usepackage{graphicx}
\usepackage{amsmath,amssymb}
\usepackage{booktabs}
\usepackage{hyperref}
\usepackage{xcolor}
\usepackage{titlesec}
\usepackage{enumitem}
\setlist{nosep}
\titlespacing*{\section}{0pt}{10pt}{4pt}
\titlespacing*{\subsection}{0pt}{8pt}{3pt}
\hypersetup{colorlinks=true,linkcolor=blue,urlcolor=blue,citecolor=blue}

\title{\textbf{Reproducing and Extending LLM-as-a-Judge:\\ MT-Bench, FastChat, and New Evaluation Data}}
\author{Uday Reddy Poli \\ Department of Computing, Macquarie University \\ COMP8240: Applications of Data Science, Session 2 2026 \\ \texttt{Project Final Report, Assessment Task \#3}}
\date{}

\begin{document}
\maketitle

\begin{abstract}
This report describes the completion of a semester-long project reproducing and extending
Zheng et al.'s LLM-as-a-judge paper, \emph{Judging LLM-as-a-Judge with MT-Bench
and Chatbot Arena} \cite{zheng2023judging}, using the authors' open-source FastChat codebase. The
original paper's headline claim -- that GPT-4 agrees with human preference judgments more than 80\%
of the time, comparably to human-human agreement -- was reproduced directly from the authors'
own released vote data, obtaining 85.9\% (turn 1) and 85.3\% (turn 2) GPT-4-vs-human agreement,
excluding ties. The original codebase was installed and run successfully, surfacing and fixing two
real, previously undocumented issues: missing transitive dependencies in the package's declared
install extras, and a category-registration bug that crashes grading on any new question category.
Sixteen new MT-Bench-style questions were constructed across two categories -- harder coding
problems and Australian local-context knowledge -- with unit-tested reference answers for all eight
coding questions, and were evaluated with an LLM-as-judge pass using Claude as a documented,
openly-stated substitute for the original paper's GPT-4 judge, since no OpenAI API access was
available. The substitute judge produced a mean score of 8.13/10 across the 16 new questions, with
real score variation (6--9) indicating the judge is discriminating between answer quality rather
than defaulting to uniform scores.
\end{abstract}

\section{Description of Paper, and Justification}

Zheng et al.'s \emph{Judging LLM-as-a-Judge with MT-Bench and Chatbot Arena}
\cite{zheng2023judging} (NeurIPS 2023 Datasets and Benchmarks track) addresses a practical problem
in evaluating large language models: human preference annotation is the gold standard for judging
open-ended chat quality, but it is slow and expensive to collect at scale. The paper proposes and
validates \emph{LLM-as-a-judge}: using a strong model (GPT-4) to score or compare model responses
in place of a human annotator. Its evaluation methodology compares GPT-4's judgments against a
large pool of crowdsourced human preference votes, on MT-Bench (their own constructed
multi-turn benchmark) and on Chatbot Arena conversations. On MT-Bench, the paper reports that GPT-4
agrees with human majority-vote preference more than 80\% of the time, a level of agreement
comparable to (and, excluding ties, exceeding) the agreement rate between two independent human
annotators.

The code for this project is lm-sys's FastChat repository on GitHub, specifically its
\texttt{llm\_judge} module.\footnote{\url{https://github.com/lm-sys/FastChat}} FastChat was chosen
for three reasons. First, feasibility was verified upfront rather than assumed: the repository is
public, actively maintained, and well-documented, and because the judging step is pure API calls
rather than local model inference, it does not require a GPU -- the only real requirement is an
OpenAI API key and a modest amount of paid credit. Second, the paper is highly influential: it was
published at NeurIPS 2023 (a CORE A* venue) and has accumulated over 600 citations on Google
Scholar as of this project. Third, the topic -- LLM-as-a-judge and model evaluation methodology --
sits squarely in an area directly relevant to applied data science practice, where it is
increasingly common to use one model to evaluate another's outputs cheaply.

\section{Description of the Original Dataset}

The paper's benchmark is MT-Bench: 80 carefully designed multi-turn questions spanning eight
categories -- writing, roleplay, extraction, reasoning, math, coding, STEM, and humanities/social
science -- with ten questions per category. Each question has exactly two turns, so a model being
evaluated must hold context across a follow-up instruction rather than answer a single isolated
prompt. For two categories, math and coding, the dataset also ships a reference answer per question,
used to ground the judge's grading in those categories (FastChat's \texttt{NEED\_REF\_CATS} list).
Figure~\ref{fig:question-snippet} shows a verified snippet confirming the dataset loads correctly.

\begin{figure}[h]
\centering
\footnotesize
\begin{verbatim}
$ python -c "import json; data=[
  json.loads(l) for l in
  open('data/mt_bench/question.jsonl')];
  print(len(data),'questions');
  print(sorted(set(
    d['category'] for d in data)))"
80 questions
['coding', 'extraction', 'humanities',
 'math', 'reasoning', 'roleplay',
 'stem', 'writing']
\end{verbatim}
\caption{Verified output confirming all 80 MT-Bench questions and 8 categories load correctly from
the original dataset file.}
\label{fig:question-snippet}
\end{figure}

The dataset is distributed directly inside the FastChat repository
(\texttt{fastchat/llm\_judge/data/mt\_bench/question.jsonl}) in a simple JSON-lines schema: each
line has a \texttt{question\_id}, a \texttt{category}, and a \texttt{turns} list of the two prompt
strings. Reference answers, where present, follow the same per-question-id JSON-lines format in a
separate \texttt{reference\_answer.jsonl} file. The authors additionally released the original
GPT-4 judge votes and 3.3K human preference votes used to produce the paper's headline agreement
figure, as a public dataset on Hugging Face
(\texttt{lmsys/mt\_bench\_human\_judgments}),\footnote{\url{https://huggingface.co/datasets/lmsys/mt_bench_human_judgments}}
in exactly the vote-record format FastChat's own \texttt{compute\_agreement.py} script expects --
this proved essential for Section~\ref{sec:replication} below.

\section{Replication of Original Work}
\label{sec:replication}

\subsection{Environment and running the original code}

The original pipeline has two stages. \texttt{gen\_api\_answer.py} calls each candidate model's API
and saves its answers to the 80 MT-Bench questions. \texttt{gen\_judgment.py} then sends those
answers -- plus the question, and for math/coding a reference answer -- to GPT-4, which returns a
numeric judgment per answer. The environment was set up and verified hands-on, twice: once in a
cloud working environment, and independently again on a personal MacBook Air (Python 3.12.7,
venv), to rule out an environment-specific fluke. Both times, FastChat was cloned, its
\texttt{llm\_judge} dependencies installed, and all 80 questions across all 8 categories were
confirmed to load and the scripts to execute without errors (Figure~\ref{fig:question-snippet}).
All project scripts, patches, and test files described in this report are available in the
project's public GitHub repository.\footnote{\url{https://github.com/REPLACE_WITH_USERNAME/comp8240-llm-judge-project} --
\textbf{see note to self in the submitted README: replace this footnote URL with the actual repo
link before final PDF compile.}}

\subsection{Two real issues found and fixed}

Installing and actually importing the judging code (not just reading it) surfaced two genuine,
previously undocumented problems, beyond the two already-known install quirks from the proposal
stage (a \texttt{wavedrom} build failure, and an \texttt{openai<1} version pin needed for API
compatibility).

\textbf{Undeclared transitive dependencies.} The package's own declared \texttt{llm\_judge} install
extras (in \texttt{pyproject.toml}) list only \texttt{openai<1}, \texttt{anthropic>=0.3}, and
\texttt{ray}. Actually importing \texttt{fastchat/llm\_judge/common.py} end-to-end surfaced three
more hard requirements not visible from the declared extras: \texttt{anthropic} is imported
unconditionally even when not used as a judge, and, transitively, \texttt{torch},
\texttt{transformers}, and \texttt{accelerate} are required because \texttt{common.py}
unconditionally imports \texttt{fastchat.model.model\_adapter} -- a module that is actually only
needed for \emph{local} model inference, not pure API-based judging. All four were installed and
the import now succeeds cleanly.

\textbf{Category-registration bug.} Adding a new question category and attempting to grade it
raises \texttt{KeyError: 'au\_context'} from the \texttt{temperature\_config} dictionary in
\texttt{common.py}, because that dictionary has no default or fallback entry -- every category used
for grading must be registered by name in advance. This was fixed with a one-line patch
(\texttt{"au\_context": 0.7,}, matching the temperature already used for the similarly open-ended
\texttt{writing}/\texttt{roleplay} categories, since \texttt{au\_context} questions are not
reference-checked). The fix was verified: the KeyError no longer occurs, and the category receives
a sensible non-deterministic sampling temperature.

\subsection{Headline result, reproduced}

Running a full live GPT-4 judging pass over all 80 original questions was not required to
reproduce the paper's central claim, because the authors released their own original GPT-4 votes
and 3.3K human votes as a public dataset in exactly the format their own agreement script expects
(Section 2). This made it possible to reproduce the headline result directly and for free, with
zero live API spend, using \texttt{reproduce\_headline\_agreement.py}, which loads that dataset and
computes both GPT-4-vs-human and human-vs-human agreement, with and without ties, mirroring the
original script's logic. (Hugging Face is blocked by the cloud sandbox's network policy, so this
script was instead run on a personal machine.) Table~\ref{tab:headline} shows the result.

\begin{table}[h]
\centering
\small
\begin{tabular}{@{}lll@{}}
\toprule
\textbf{Turn} & \textbf{GPT-4 vs human} & \textbf{human vs human} \\
\midrule
1 & 85.9\% (911/1060) & 82.0\% (633/772) \\
2 & 85.3\% (912/1069) & 83.4\% (649/778) \\
\bottomrule
\end{tabular}
\caption{Reproduced agreement rates, excluding ties, from the authors' own released vote data.}
\label{tab:headline}
\end{table}

Both turns reproduce the paper's ``$>$80\%'' headline claim, and in both turns GPT-4's agreement
with human preference is actually slightly \emph{higher} than human-vs-human agreement itself,
excluding ties -- consistent with the paper's own finding that GPT-4 agreement is comparable to,
or exceeds, the level of agreement between two independent human annotators. This is the strongest
piece of evidence for replication in this project: a real number, independently reproduced from the
authors' own data, not merely a claim that the code ``should work.''

\section{Construction of New Data}

Two categories of new data were constructed, both following FastChat's exact two-turn question
schema (\texttt{question\_id}, \texttt{category}, \texttt{turns}), so they drop directly into the
existing pipeline without any format conversion.

\textbf{Eight harder coding questions} (ids 1001--1008): linked-list cycle detection, bounded-memory
top-$k$ frequency counting over a log file, a buggy binary search requiring a fix, a SQL
\texttt{DENSE\_RANK} window-function query, a race condition and its fix, a thread-safe token-bucket
rate limiter, bracket/tag balance checking, and closest-pair-of-points via the haversine formula.
These were chosen to be more varied and harder than the original benchmark's coding questions, and
because \texttt{coding} is in FastChat's \texttt{NEED\_REF\_CATS} list, each needs a reference
answer to be judged in reference-guided mode.

\textbf{Eight Australian local-context questions} (ids 2001--2008): small-business grants, GST vs
payroll tax, local government structure, preferential voting, the Medicare Safety Net, buying a
first home, visa pathways, and daylight saving -- a kind of local, time-sensitive factual knowledge
the original benchmark never covers.

\subsection{Reference answers and their testing}

Real, correct reference solutions were written for all 8 coding questions, and \textbf{all 8 are
unit-tested} in \texttt{test\_refs.py} -- executed against concrete inputs, not eyeballed, including
a 100-thread concurrency stress test for the rate limiter (confirming the lock caps allowed
requests at exactly the bucket's capacity despite 100 simultaneous callers) and a temp-file-backed
test for the IP-counting question. Running \texttt{python test\_refs.py} prints
\texttt{ALL 8 REFERENCE ANSWERS VERIFIED} on a clean run. Figure~\ref{fig:ref-test} shows a
representative snippet of this.

\begin{figure}[h]
\centering
\footnotesize
\begin{verbatim}
$ python test_refs.py
1003 OK
1007 part1 OK
1007 part2 OK
1001 OK
1004 OK
1002 part1 OK
1002 part2 OK
1005 OK
1006 OK
1008 OK
ALL 8 REFERENCE ANSWERS VERIFIED
\end{verbatim}
\caption{Output of the reference-answer test suite: all 8 coding reference answers pass.}
\label{fig:ref-test}
\end{figure}

One question's wording was corrected during this testing process. Question 1002's turn 1 originally
asked for a solution using ``bounded memory,'' but the reference solution's own memory use is
$O(\text{number of distinct IPs})$, not a fixed bound -- the turn-1 solution's own documentation
already acknowledged this. The question now asks to solve it ``without loading the entire file into
memory,'' which is what the turn-1 solution actually achieves; the genuinely fixed-memory,
approximate version (a count-min-sketch-style structure) is introduced in turn 2 instead. This
makes the two-turn progression internally accurate: turn 1 is exact but distinct-IP-bounded, turn 2
is approximate but truly fixed-memory.

\subsection{Annotation and quality control}

Rather than treat the 16 questions as self-evidently correct, each was reviewed against an explicit
criterion before inclusion, recorded in full in \texttt{annotation\_record.md}. Coding questions
were reviewed against \emph{executable correctness}: does the reference answer, run against concrete
test inputs, produce the behaviour the question asks for? All 8 passed this review (post-fix). The
Australian-context questions were reviewed against a different criterion, chosen deliberately
because these questions touch facts that change over time (grant amounts, thresholds, visa
categories): not ``is there one fixed correct answer,'' but ``does the question push the model
toward a defensible, well-framed answer, with appropriate uncertainty where figures could have
changed.'' This criterion was carried forward into the judging pass in Section~\ref{sec:results}.

No crowdsourced human annotation was used for this new data -- all construction and review was
carried out by a single annotator (the author), which is itself discussed as a limitation in
Section~\ref{sec:ethics}, since it means no inter-annotator agreement measure exists for the new
questions, unlike the original paper's large crowdsourced human-vote pool.

\section{Results on New Data}
\label{sec:results}

Evaluating the 16 new questions with the original pipeline requires an LLM judge, but no OpenAI API
key was available for a live GPT-4 run through \texttt{gen\_api\_answer.py} /
\texttt{gen\_judgment.py}. Rather than report a pure limitation with no results, this stage used
Claude (via the assistant session supporting this project) as a documented, free substitute for
\emph{both} roles the original pipeline needs: generating one independent set of 16 candidate
answers (\texttt{candidate\_answer\_additions.jsonl}, written without reference to the reference
answers), and then, in a separate judging pass, scoring those candidate answers on FastChat's 1--10
single-answer grading scale -- reference-guided against \texttt{reference\_answer\_additions.jsonl}
for the coding category, and against the factual-correctness-plus-appropriate-uncertainty criterion
from Section 4.2 for the Australian-context category. Full scores and per-question notes are in
\texttt{judge\_results.md}.

\textbf{This substitution is a stated methodology deviation, not a hidden one.} Using the same
underlying model family for both the candidate and the judge role carries a known self-evaluation
bias risk discussed in the LLM-as-judge literature, including in the original paper itself,
and Claude is not GPT-4, the specific model the original paper evaluates. The results below should
therefore be read as evidence that the new questions and their grading criteria function correctly
end-to-end, not as a reproduction of what a live GPT-4 judging run would score.

\begin{table}[h]
\centering
\small
\begin{tabular}{@{}lccc@{}}
\toprule
\textbf{Category} & \textbf{Turn 1} & \textbf{Turn 2} & \textbf{Overall} \\
\midrule
Coding (n=8) & 8.50 & 7.88 & 8.19 \\
Au-context (n=8) & 8.13 & 8.00 & 8.06 \\
\midrule
All (n=16) & 8.31 & 7.94 & \textbf{8.13} \\
\bottomrule
\end{tabular}
\caption{Mean judge scores (1--10 scale) on the 16 new questions.}
\label{tab:results}
\end{table}

Scores ranged from 6 to 9 across the 32 individual (question, turn) gradings -- real variation,
not a uniform ceiling score, which is useful evidence that the grading criteria are actually
discriminating between stronger and weaker answers. The single lowest-scoring answer (1002, turn 2,
score 6) was a candidate answer whose eviction-based memory-bounding strategy was functionally
adequate but noticeably cruder than the reference's count-min-sketch approach -- exactly the kind of
difference a reference-guided judge should catch.

What this does \emph{not} show is GPT-4-vs-human agreement specifically on this new data, since
no human annotation pass or live GPT-4 run was performed on these 16 questions, unlike the original
80-question result reproduced in Section~\ref{sec:replication} from the authors' own released human
votes. Running a live \texttt{gen\_api\_answer.py} / \texttt{gen\_judgment.py} pass against GPT-4 on
this new data, budget permitting, remains the natural next step and is discussed further in
Section~\ref{sec:conclusion}.

\section{Ethical Considerations}
\label{sec:ethics}

\textbf{In the original paper.} The original paper's evaluation relies on a large pool of
crowdsourced human preference annotators. This raises standard crowdsourcing ethics questions the
paper itself discusses only briefly: annotator compensation and working conditions, the
demographic composition of the annotator pool relative to the population whose preferences are
being modelled, and the risk that preference judgments reflect superficial qualities (length,
confident tone, formatting) rather than genuine correctness or helpfulness -- a bias the
LLM-as-judge literature has since documented further. Using a single commercial model (GPT-4) as
judge also concentrates evaluation power in one proprietary system controlled by one company,
raising fairness and access concerns: researchers without budget for GPT-4 API access cannot fully
replicate or audit the judging step, which is itself directly illustrated by this project's own
budget constraint in Section~\ref{sec:results}.

\textbf{For this project's new data.} If the Australian-context questions and their model
evaluations were to be published as a research dataset, under the Australian Code for the
Responsible Conduct of Research \cite{aust_code_2018}, the most relevant considerations are: (1)
\emph{data currency and potential for misinformation} -- several questions touch rules that change
over time (grant thresholds, visa categories, Medicare Safety Net amounts), so any published dataset
or model trained/evaluated against it would need a clear ``as-of'' date and explicit framing that
answers should not be treated as current legal or financial advice; (2) \emph{single-annotator bias}
-- all 16 questions, their review criteria, and (in this report) the substitute judge's scores were
produced by one person/model pipeline, with no inter-annotator agreement measure, unlike the
original paper's large human-vote pool; publishing this data as a benchmark would require either
a second independent annotator or explicit disclosure of this limitation; (3) \emph{no personal or
sensitive data is involved} -- all 16 questions concern public government policy information or
general programming problems, not individuals, so there are no participant-consent or privacy
concerns of the kind the Code's human-research provisions address; and (4) \emph{transparency about
the substitute judge}, consistent with the Code's research-integrity principles around honest
reporting of methods and limitations -- which is why Section~\ref{sec:results} states the
GPT-4-to-Claude substitution explicitly rather than presenting the scores as equivalent to the
original paper's judging methodology.

\section{Conclusions, and Other Reflections}
\label{sec:conclusion}

This project set out to replicate Zheng et al.'s LLM-as-a-judge paper and extend
it to new data. Replication succeeded on every piece attempted: the original codebase installs and
runs cleanly (after fixing two real, previously undocumented issues), and the paper's headline
$>$80\% GPT-4-vs-human agreement claim was reproduced directly from the authors' own released vote
data (85.9\% and 85.3\% across the two turns, excluding ties), exceeding human-vs-human agreement
itself. Extension to new data was also completed: 16 new two-turn questions were constructed across
two categories not covered by the original benchmark, with all 8 coding reference answers written
and unit-tested, and a documented, openly-stated substitute LLM-as-judge pass was run across all 16
questions, producing real, non-uniform scores (mean 8.13/10, range 6--9) that demonstrate the
questions and grading criteria work end-to-end.

The main limitation, acknowledged throughout rather than hidden, is budget: no OpenAI API access
was available, so neither the original 80 questions nor the 16 new questions could be run through a
live GPT-4 judging pass directly. The 80-question result was still reproduced exactly, because the
authors generously released their own vote data; the 16-question result could not be reproduced the
same way (it is new data with no existing released votes), so a substitute judge was used instead,
with its limitations stated plainly in Section~\ref{sec:results}. If given further time and budget,
the clearest next step would be running \texttt{gen\_api\_answer.py} and \texttt{gen\_judgment.py}
against GPT-4 directly on the 16 new questions, and ideally collecting a small human-preference
pass (even just 2--3 annotators) on a subset, to get a genuine agreement-rate figure for the new
data comparable to Table~\ref{tab:headline}'s figure for the original 80 questions. A second
reflection is on annotation process: constructing the Australian-context questions surfaced how much
harder it is to write a good open-ended factual question with a \emph{stable} grading criterion than
to write one with a single fixed answer -- several early draft questions (not included here) asked
for specific dollar figures that would have gone stale within a year, which is what motivated the
``factual correctness of stable structure plus appropriate uncertainty'' criterion used throughout
Section 4.2, rather than exact-answer matching.

\bibliographystyle{plain}
\begin{thebibliography}{9}

\bibitem[Zheng et al., 2023]{zheng2023judging}
Lianmin Zheng, Wei-Lin Chiang, Ying Sheng, Siyuan Zhuang, Zhanghao Wu, Yonghao Zhuang, Zi Lin,
Zhuohan Li, Dacheng Li, Eric P. Xing, Hao Zhang, Joseph E. Gonzalez, and Ion Stoica. 2023.
Judging LLM-as-a-Judge with MT-Bench and Chatbot Arena. In \emph{Advances in Neural Information
Processing Systems (NeurIPS), Datasets and Benchmarks Track}.
\url{https://github.com/lm-sys/FastChat}

\bibitem[Chen et al., 2021]{chen2021humaneval}
Mark Chen, Jerry Tworek, Heewoo Jun, Qiming Yuan, et al. 2021. Evaluating Large Language Models
Trained on Code. \emph{arXiv preprint arXiv:2107.03374}.

\bibitem[Australian Code, 2018]{aust_code_2018}
National Health and Medical Research Council, Australian Research Council, and Universities
Australia. 2018. \emph{Australian Code for the Responsible Conduct of Research}, 2018 edition.
Canberra: National Health and Medical Research Council.

\end{thebibliography}

\end{document}
